\section{The Research Journey}\label{sec:journey}

FOA was not our first design. Figure~\ref{fig:journey} traces how each idea answered the weakness of the
one before, and what each step taught us.

\begin{figure*}[t]
\centering
\begin{tikzpicture}[node distance=0.32cm,
  st/.style={draw=black!45, rounded corners=2pt, fill=black!4, text width=2.0cm, minimum height=1.55cm, align=center, font=\scriptsize, inner sep=3pt},
  hi/.style={st, draw=blue!70!black, fill=blue!6},
  ar/.style={-{Stealth[length=2mm]}, black!55, thick}]
\node[st] (a) {\textbf{1. Lone device}\\ learns its own clean traffic\\[2pt] \textit{poisoned in 39\% of runs}};
\node[st, right=of a] (b) {\textbf{2. Stop when suspicious}\\ gates and switch-offs\\[2pt] \textit{frozen by decoys: 100\% false alarms}};
\node[st, right=of b] (c) {\textbf{3. TOA}\\ refuse moves toward the attack\\[2pt] \textit{5.5\%, but needs the attack in advance}};
\node[hi, right=of c] (d) {\textbf{4. FOA}\\ neighbours as witnesses, budgeted rest\\[2pt] \textit{10.8\%, no attack knowledge}};
\node[st, right=of d] (e) {\textbf{5. Real boards}\\ normal region collapsed\\[2pt] \textit{under 2\% of its size, 90\% alarms}};
\node[hi, right=of e] (f) {\textbf{6. Coverage floor}\\ normal never shrinks below calibration\\[2pt] \textit{held at 1.0 on every board}};
\node[hi, right=of f] (g) {\textbf{7. Calibration fix}\\ uneven sampling gaps\\[2pt] \textit{board false alarms 25\% to 1.6\%}};
\foreach \x/\y in {a/b,b/c,c/d,d/e,e/f,f/g} \draw[ar] (\x) -- (\y);
\end{tikzpicture}
\caption{How the design evolved. Each step fixed the weakness the previous one exposed.}\label{fig:journey}
\end{figure*}

\textbf{From a lone device to gates.} Self-labelled learning on one device is the natural starting point,
and it is exactly what the boiling frog exploits. Stopping learning on suspicion protects against it but
introduces the freeze attack (Section~\ref{sec:eval}): five percent decoys locked the alarm switch-off at
\KLFrozenFPR\ false alarms.

\textbf{From TOA to FOA.} Our first method, Threat-Orthogonal Adaptation, refused movement toward the
known attack direction. It is very safe (5.5\% compromised) but needs the attack's signature in advance,
which a deployed device never has, and it raised 38\% false alarms. FOA replaces that oracle knowledge
with a witness a device does have: its neighbours.

\textbf{From simulation to boards.} On three ESP8266 boards the detectors shrank their normal regions in
agreement until almost everything alarmed (Section~\ref{sec:hw}). The cause, a guard on growth with none
on shrinking, was reproduced in simulation with no attacker. The coverage floor fixed it, held on every
board, and corrected our own claim: FOA's lasting advantage is safety under drift. A second fault,
calibration sampling in step with the 200\,ms telemetry, explained the remaining false alarms; uneven
sampling gaps brought them to 1.5\% to 1.7\%.

\textbf{Lessons.} (i) A defence that can stop learning can be made to stop. (ii) A second witness beats
a stronger single witness. (iii) Guard both directions of change, not only the one an attacker uses.
(iv) Test on hardware early: the boards found two faults that thousands of simulated runs did not.
